Proposition~\ref{ms:prop-2-2} identifies $m_0$ as the multiplier symbol. Since $|\xi|^2-\xi_b^2=|\xi_x|^2$ and $|\xi|+\xi_b=\sqrt2\kappa$,
\[
 |\xi|=\kappa(1+|v|^2/2)/\sqrt2 .
\]
The annulus bounds and $|v|\le1/10$ give $\kappa\asymp h^{-1}$. The bound $|\operatorname{Re}(W)\cdot v|\le1/100$ gives $b_0\asymp1$. Inverting $m_0$ gives the claimed reciprocal. If $y\asymp h$, then $\Psi$ is bounded and $(\kappa/y)^{d/2}\asymp h^{-d}$, giving the two-sided bound.

In the original torus frame the frequencies belong to $(2\pi/100)\mathbb Z^n$. The radius-$C_0/h$ ball about the origin contains the annulus. Pairwise disjoint lattice cubes centered on points of the ball, of side $2\pi/100$, fit inside a ball of radius $C_0/h+\pi\sqrt n/100$. Volume comparison bounds their number by $Ch^{-n}$, and rotation preserves the radii and volumes. Finally Parseval gives $\sum_\xi|c_\xi|^2=100^{-n}\|f\|_2^2$, so Cauchy--Schwarz yields the coefficient bound.
